%%%%%%%%%%%%%%%%%%%%%%%%%%%%%%%%%%%%%%%%%
% The Legrand Orange Book
% LaTeX Template
% Version 2.1 (14/11/15)
%
% This template has been downloaded from:
% http://www.LaTeXTemplates.com
%
% Mathias Legrand (legrand.mathias@gmail.com) with modifications by:
% Vel (vel@latextemplates.com)
%
% License:
% CC BY-NC-SA 3.0 (http://creativecommons.org/licenses/by-nc-sa/3.0/)
%
% Compiling this template:
% This template uses biber for its bibliography and makeindex for its index.
% When you first open the template, compile it from the command line with the 
% commands below to make sure your LaTeX distribution is configured correctly:
%
% 1) pdflatex main
% 2) makeindex main.idx -s StyleInd.ist
% 3) biber main
% 4) pdflatex main x 2
%
% After this, when you wish to update the bibliography/index use the appropriate
% command above and make sure to compile with pdflatex several times 
% afterwards to propagate your changes to the document.
%
% This template also uses a number of packages which may need to be
% updated to the newest versions for the template to compile. It is strongly
% recommended you update your LaTeX distribution if you have any
% compilation errors.
%
% Important note:
% Chapter heading images should have a 2:1 width:height ratio,
% e.g. 920px width and 460px height.
%
%%%%%%%%%%%%%%%%%%%%%%%%%%%%%%%%%%%%%%%%%

%----------------------------------------------------------------------------------------
%	PACKAGES AND OTHER DOCUMENT CONFIGURATIONS
%----------------------------------------------------------------------------------------

\documentclass[11pt,fleqn]{book} % Default font size and left-justified equations

\newcommand{\TBC} {\vspace {1em} \noindent [SAMPLE TEXT. TO BE COMPLETED IN THE FINAL VERSION.] \vspace {1em}}

%----------------------------------------------------------------------------------------

%%%%%%%%%%%%%%%%%%%%%%%%%%%%%%%%%%%%%%%%%
% The Legrand Orange Book
% Structural Definitions File
% Version 2.0 (9/2/15)
%
% Original author:
% Mathias Legrand (legrand.mathias@gmail.com) with modifications by:
% Vel (vel@latextemplates.com)
% 
% This file has been downloaded from:
% http://www.LaTeXTemplates.com
%
% License:
% CC BY-NC-SA 3.0 (http://creativecommons.org/licenses/by-nc-sa/3.0/)
%
%%%%%%%%%%%%%%%%%%%%%%%%%%%%%%%%%%%%%%%%%

%----------------------------------------------------------------------------------------
%	VARIOUS REQUIRED PACKAGES AND CONFIGURATIONS
%----------------------------------------------------------------------------------------

\usepackage[top=3cm,bottom=3cm,left=3cm,right=3cm,headsep=10pt,a4paper]{geometry} % Page margins

\usepackage{upgreek} 
\usepackage{graphicx} % Required for including pictures
\graphicspath{{Pictures/}} % Specifies the directory where pictures are stored

\usepackage{lipsum} % Inserts dummy text

\usepackage{tikz} % Required for drawing custom shapes

\usepackage[english]{babel} % English language/hyphenation

\usepackage{enumitem} % Customize lists
\setlist{nolistsep} % Reduce spacing between bullet points and numbered lists

\usepackage{booktabs} % Required for nicer horizontal rules in tables

\usepackage{xcolor} % Required for specifying colors by name
% \definecolor{ocre}{RGB}{243,102,25} % Define the orange color used for highlighting throughout the book
\definecolor{ocre}{RGB}{203,142,95} % Define the orange color used for highlighting throughout the book

%----------------------------------------------------------------------------------------
%	FONTS
%----------------------------------------------------------------------------------------

\usepackage{avant} % Use the Avantgarde font for headings
%\usepackage{times} % Use the Times font for headings
\usepackage{mathptmx} % Use the Adobe Times Roman as the default text font together with math symbols from the Sym­bol, Chancery and Com­puter Modern fonts

\usepackage{microtype} % Slightly tweak font spacing for aesthetics
\usepackage[utf8]{inputenc} % Required for including letters with accents
\usepackage[T1]{fontenc} % Use 8-bit encoding that has 256 glyphs

%----------------------------------------------------------------------------------------
%	BIBLIOGRAPHY AND INDEX
%----------------------------------------------------------------------------------------
% \usepackage[style=alphabetic,citestyle=numeric,sorting=nyt,sortcites=true,autopunct=true,babel=hyphen,hyperref=true,abbreviate=false,backref=true,backend=biber]{biblatex}
% \addbibresource{bibliography.bib} % BibTeX bibliography file
% \defbibheading{bibempty}{}
% \usepackage[pdftex]{hyperref}


\usepackage{calc} % For simpler calculation - used for spacing the index letter headings correctly
\usepackage{makeidx} % Required to make an index
\makeindex % Tells LaTeX to create the files required for indexing

%----------------------------------------------------------------------------------------
%	MAIN TABLE OF CONTENTS
%----------------------------------------------------------------------------------------

\usepackage{titletoc} % Required for manipulating the table of contents

\contentsmargin{0cm} % Removes the default margin

% Part text styling
\titlecontents{part}[0cm]
{\addvspace{20pt}\centering\large\bfseries}
{}
{}
{}

% Chapter text styling
\titlecontents{chapter}[1.25cm] % Indentation
{\addvspace{12pt}\large\sffamily\bfseries} % Spacing and font options for chapters
{\color{ocre!60}\contentslabel[\Large\thecontentslabel]{1.25cm}\color{ocre}} % Chapter number
{\color{ocre}}  
{\color{ocre!60}\normalsize\;\titlerule*[.5pc]{.}\;\thecontentspage} % Page number

% Section text styling
\titlecontents{section}[1.25cm] % Indentation
{\addvspace{3pt}\sffamily\bfseries} % Spacing and font options for sections
{\contentslabel[\thecontentslabel]{1.25cm}} % Section number
{}
{\hfill\color{black}\thecontentspage} % Page number
[]

% Subsection text styling
\titlecontents{subsection}[1.25cm] % Indentation
{\addvspace{1pt}\sffamily\small} % Spacing and font options for subsections
{\contentslabel[\thecontentslabel]{1.25cm}} % Subsection number
{}
{\ \titlerule*[.5pc]{.}\;\thecontentspage} % Page number
[]

% List of figures
\titlecontents{figure}[0em]
{\addvspace{-5pt}\sffamily}
{\thecontentslabel\hspace*{1em}}
{}
{\ \titlerule*[.5pc]{.}\;\thecontentspage}
[]

% List of tables
\titlecontents{table}[0em]
{\addvspace{-5pt}\sffamily}
{\thecontentslabel\hspace*{1em}}
{}
{\ \titlerule*[.5pc]{.}\;\thecontentspage}
[]

%----------------------------------------------------------------------------------------
%	MINI TABLE OF CONTENTS IN PART HEADS
%----------------------------------------------------------------------------------------

% Chapter text styling
\titlecontents{lchapter}[0em] % Indenting
{\addvspace{15pt}\large\sffamily\bfseries} % Spacing and font options for chapters
{\color{ocre}\contentslabel[\Large\thecontentslabel]{1.25cm}\color{ocre}} % Chapter number
{}  
{\color{ocre}\normalsize\sffamily\bfseries\;\titlerule*[.5pc]{.}\;\thecontentspage} % Page number

% Section text styling
\titlecontents{lsection}[0em] % Indenting
{\sffamily\small} % Spacing and font options for sections
{\contentslabel[\thecontentslabel]{1.25cm}} % Section number
{}
{}

% Subsection text styling
\titlecontents{lsubsection}[.5em] % Indentation
{\normalfont\footnotesize\sffamily} % Font settings
{}
{}
{}

%----------------------------------------------------------------------------------------
%	PAGE HEADERS
%----------------------------------------------------------------------------------------

\usepackage{fancyhdr} % Required for header and footer configuration

\pagestyle{fancy}
\renewcommand{\chaptermark}[1]{\markboth{\sffamily\normalsize\bfseries\chaptername\ \thechapter.\ #1}{}} % Chapter text font settings
\renewcommand{\sectionmark}[1]{\markright{\sffamily\normalsize\thesection\hspace{5pt}#1}{}} % Section text font settings
\fancyhf{} \fancyhead[LE,RO]{\sffamily\normalsize\thepage} % Font setting for the page number in the header
\fancyhead[LO]{\rightmark} % Print the nearest section name on the left side of odd pages
\fancyhead[RE]{\leftmark} % Print the current chapter name on the right side of even pages
\renewcommand{\headrulewidth}{0.5pt} % Width of the rule under the header
\addtolength{\headheight}{2.5pt} % Increase the spacing around the header slightly
\renewcommand{\footrulewidth}{0pt} % Removes the rule in the footer
\fancypagestyle{plain}{\fancyhead{}\renewcommand{\headrulewidth}{0pt}} % Style for when a plain pagestyle is specified

% Removes the header from odd empty pages at the end of chapters
\makeatletter
\renewcommand{\cleardoublepage}{
\clearpage\ifodd\c@page\else
\hbox{}
\vspace*{\fill}
\thispagestyle{empty}
\newpage
\fi}

%----------------------------------------------------------------------------------------
%	THEOREM STYLES
%----------------------------------------------------------------------------------------

\usepackage{amsmath,amsfonts,amssymb,amsthm} % For math equations, theorems, symbols, etc

\newcommand{\intoo}[2]{\mathopen{]}#1\,;#2\mathclose{[}}
\newcommand{\ud}{\mathop{\mathrm{{}d}}\mathopen{}}
\newcommand{\intff}[2]{\mathopen{[}#1\,;#2\mathclose{]}}
\newtheorem{notation}{Notation}[chapter]

% Boxed/framed environments
\newtheoremstyle{ocrenumbox}% % Theorem style name
{0pt}% Space above
{0pt}% Space below
{\normalfont}% % Body font
{}% Indent amount
{\small\bf\sffamily\color{ocre}}% % Theorem head font
{\;}% Punctuation after theorem head
{0.25em}% Space after theorem head
{\small\sffamily\color{ocre}\thmname{#1}\nobreakspace\thmnumber{\@ifnotempty{#1}{}\@upn{#2}}% Theorem text (e.g. Theorem 2.1)
\thmnote{\nobreakspace\the\thm@notefont\sffamily\bfseries\color{black}---\nobreakspace#3.}} % Optional theorem note
\renewcommand{\qedsymbol}{$\blacksquare$}% Optional qed square

\newtheoremstyle{blacknumex}% Theorem style name
{5pt}% Space above
{5pt}% Space below
{\normalfont}% Body font
{} % Indent amount
{\small\bf\sffamily}% Theorem head font
{\;}% Punctuation after theorem head
{0.25em}% Space after theorem head
{\small\sffamily{\tiny\ensuremath{\blacksquare}}\nobreakspace\thmname{#1}\nobreakspace\thmnumber{\@ifnotempty{#1}{}\@upn{#2}}% Theorem text (e.g. Theorem 2.1)
\thmnote{\nobreakspace\the\thm@notefont\sffamily\bfseries---\nobreakspace#3.}}% Optional theorem note

\newtheoremstyle{blacknumbox} % Theorem style name
{0pt}% Space above
{0pt}% Space below
{\normalfont}% Body font
{}% Indent amount
{\small\bf\sffamily}% Theorem head font
{\;}% Punctuation after theorem head
{0.25em}% Space after theorem head
{\small\sffamily\thmname{#1}\nobreakspace\thmnumber{\@ifnotempty{#1}{}\@upn{#2}}% Theorem text (e.g. Theorem 2.1)
\thmnote{\nobreakspace\the\thm@notefont\sffamily\bfseries---\nobreakspace#3.}}% Optional theorem note

% Non-boxed/non-framed environments
\newtheoremstyle{ocrenum}% % Theorem style name
{5pt}% Space above
{5pt}% Space below
{\normalfont}% % Body font
{}% Indent amount
{\small\bf\sffamily\color{ocre}}% % Theorem head font
{\;}% Punctuation after theorem head
{0.25em}% Space after theorem head
{\small\sffamily\color{ocre}\thmname{#1}\nobreakspace\thmnumber{\@ifnotempty{#1}{}\@upn{#2}}% Theorem text (e.g. Theorem 2.1)
\thmnote{\nobreakspace\the\thm@notefont\sffamily\bfseries\color{black}---\nobreakspace#3.}} % Optional theorem note
\renewcommand{\qedsymbol}{$\blacksquare$}% Optional qed square
\makeatother

% Defines the theorem text style for each type of theorem to one of the three styles above
\newcounter{dummy} 
\numberwithin{dummy}{section}
\theoremstyle{ocrenumbox}
\newtheorem{theoremeT}[dummy]{Theorem}
\newtheorem{problem}{Problem}[chapter]
\newtheorem{exerciseT}{Exercise}[chapter]
\theoremstyle{blacknumex}
\newtheorem{exampleT}{Example}[chapter]
\theoremstyle{blacknumbox}
\newtheorem{vocabulary}{Vocabulary}[chapter]
\newtheorem{definitionT}{Definition}[section]
\newtheorem{corollaryT}[dummy]{Corollary}
\theoremstyle{ocrenum}
\newtheorem{proposition}[dummy]{Proposition}

%----------------------------------------------------------------------------------------
%	DEFINITION OF COLORED BOXES
%----------------------------------------------------------------------------------------

\RequirePackage[framemethod=default]{mdframed} % Required for creating the theorem, definition, exercise and corollary boxes

% Theorem box
\newmdenv[skipabove=7pt,
skipbelow=7pt,
backgroundcolor=black!5,
linecolor=ocre,
innerleftmargin=5pt,
innerrightmargin=5pt,
innertopmargin=5pt,
leftmargin=0cm,
rightmargin=0cm,
innerbottommargin=5pt]{tBox}

% Exercise box	  
\newmdenv[skipabove=7pt,
skipbelow=7pt,
rightline=false,
leftline=true,
topline=false,
bottomline=false,
backgroundcolor=ocre!10,
linecolor=ocre,
innerleftmargin=5pt,
innerrightmargin=5pt,
innertopmargin=5pt,
innerbottommargin=5pt,
leftmargin=0cm,
rightmargin=0cm,
linewidth=4pt]{eBox}	

% Definition box
\newmdenv[skipabove=7pt,
skipbelow=7pt,
rightline=false,
leftline=true,
topline=false,
bottomline=false,
linecolor=ocre,
innerleftmargin=5pt,
innerrightmargin=5pt,
innertopmargin=0pt,
leftmargin=0cm,
rightmargin=0cm,
linewidth=4pt,
innerbottommargin=0pt]{dBox}	

% Corollary box
\newmdenv[skipabove=7pt,
skipbelow=7pt,
rightline=false,
leftline=true,
topline=false,
bottomline=false,
linecolor=gray,
backgroundcolor=black!5,
innerleftmargin=5pt,
innerrightmargin=5pt,
innertopmargin=5pt,
leftmargin=0cm,
rightmargin=0cm,
linewidth=4pt,
innerbottommargin=5pt]{cBox}

% Creates an environment for each type of theorem and assigns it a theorem text style from the "Theorem Styles" section above and a colored box from above
\newenvironment{theorem}{\begin{tBox}\begin{theoremeT}}{\end{theoremeT}\end{tBox}}
\newenvironment{exercise}{\begin{eBox}\begin{exerciseT}}{\hfill{\color{ocre}\tiny\ensuremath{\blacksquare}}\end{exerciseT}\end{eBox}}				  
\newenvironment{definition}{\begin{dBox}\begin{definitionT}}{\end{definitionT}\end{dBox}}	
\newenvironment{example}{\begin{exampleT}}{\hfill{\tiny\ensuremath{\blacksquare}}\end{exampleT}}		
\newenvironment{corollary}{\begin{cBox}\begin{corollaryT}}{\end{corollaryT}\end{cBox}}	

%----------------------------------------------------------------------------------------
%	REMARK ENVIRONMENT
%----------------------------------------------------------------------------------------

\newenvironment{remark}{\par\vspace{10pt}\small % Vertical white space above the remark and smaller font size
\begin{list}{}{
\leftmargin=35pt % Indentation on the left
\rightmargin=25pt}\item\ignorespaces % Indentation on the right
\makebox[-2.5pt]{\begin{tikzpicture}[overlay]
\node[draw=ocre!60,line width=1pt,circle,fill=ocre!25,font=\sffamily\bfseries,inner sep=2pt,outer sep=0pt] at (-15pt,0pt){\textcolor{ocre}{R}};\end{tikzpicture}} % Orange R in a circle
\advance\baselineskip -1pt}{\end{list}\vskip5pt} % Tighter line spacing and white space after remark

%----------------------------------------------------------------------------------------
%	SECTION NUMBERING IN THE MARGIN
%----------------------------------------------------------------------------------------

\makeatletter
\renewcommand{\@seccntformat}[1]{\llap{\textcolor{ocre}{\csname the#1\endcsname}\hspace{1em}}}                    
\renewcommand{\section}{\@startsection{section}{1}{\z@}
{-4ex \@plus -1ex \@minus -.4ex}
{1ex \@plus.2ex }
{\normalfont\large\sffamily\bfseries}}
\renewcommand{\subsection}{\@startsection {subsection}{2}{\z@}
{-3ex \@plus -0.1ex \@minus -.4ex}
{0.5ex \@plus.2ex }
{\normalfont\sffamily\bfseries}}
\renewcommand{\subsubsection}{\@startsection {subsubsection}{3}{\z@}
{-2ex \@plus -0.1ex \@minus -.2ex}
{.2ex \@plus.2ex }
{\normalfont\small\sffamily\bfseries}}                        
\renewcommand\paragraph{\@startsection{paragraph}{4}{\z@}
{-2ex \@plus-.2ex \@minus .2ex}
{.1ex}
{\normalfont\small\sffamily\bfseries}}

%----------------------------------------------------------------------------------------
%	PART HEADINGS
%----------------------------------------------------------------------------------------

% numbered part in the table of contents
\newcommand{\@mypartnumtocformat}[2]{%
\setlength\fboxsep{0pt}%
\noindent\colorbox{ocre!20}{\strut\parbox[c][.7cm]{\ecart}{\color{ocre!70}\Large\sffamily\bfseries\centering#1}}\hskip\esp\colorbox{ocre!40}{\strut\parbox[c][.7cm]{\linewidth-\ecart-\esp}{\Large\sffamily\centering#2}}}%
%%%%%%%%%%%%%%%%%%%%%%%%%%%%%%%%%%
% unnumbered part in the table of contents
\newcommand{\@myparttocformat}[1]{%
\setlength\fboxsep{0pt}%
\noindent\colorbox{ocre!40}{\strut\parbox[c][.7cm]{\linewidth}{\Large\sffamily\centering#1}}}%
%%%%%%%%%%%%%%%%%%%%%%%%%%%%%%%%%%
\newlength\esp
\setlength\esp{4pt}
\newlength\ecart
\setlength\ecart{1.2cm-\esp}
\newcommand{\thepartimage}{}%
\newcommand{\partimage}[1]{\renewcommand{\thepartimage}{#1}}%
\def\@part[#1]#2{%
\ifnum \c@secnumdepth >-2\relax%
\refstepcounter{part}%
\addcontentsline{toc}{part}{\texorpdfstring{\protect\@mypartnumtocformat{\thepart}{#1}}{\partname~\thepart\ ---\ #1}}
\else%
\addcontentsline{toc}{part}{\texorpdfstring{\protect\@myparttocformat{#1}}{#1}}%
\fi%
\startcontents%
\markboth{}{}%
{\thispagestyle{empty}%
\begin{tikzpicture}[remember picture,overlay]%
\node at (current page.north west){\begin{tikzpicture}[remember picture,overlay]%	
\fill[ocre!20](0cm,0cm) rectangle (\paperwidth,-\paperheight);
\node[anchor=north] at (4cm,-3.25cm){\color{ocre!40}\fontsize{220}{100}\sffamily\bfseries\@Roman\c@part}; 
\node[anchor=south east] at (\paperwidth-1cm,-\paperheight+1cm){\parbox[t][][t]{8.5cm}{
\printcontents{l}{0}{\setcounter{tocdepth}{1}}%
}};
\node[anchor=north east] at (\paperwidth-1.5cm,-3.25cm){\parbox[t][][t]{15cm}{\strut\raggedleft\color{white}\fontsize{30}{30}\sffamily\bfseries#2}};
\end{tikzpicture}};
\end{tikzpicture}}%
\@endpart}
\def\@spart#1{%
\startcontents%
\phantomsection
{\thispagestyle{empty}%
\begin{tikzpicture}[remember picture,overlay]%
\node at (current page.north west){\begin{tikzpicture}[remember picture,overlay]%	
\fill[ocre!20](0cm,0cm) rectangle (\paperwidth,-\paperheight);
\node[anchor=north east] at (\paperwidth-1.5cm,-3.25cm){\parbox[t][][t]{15cm}{\strut\raggedleft\color{white}\fontsize{30}{30}\sffamily\bfseries#1}};
\end{tikzpicture}};
\end{tikzpicture}}
\addcontentsline{toc}{part}{\texorpdfstring{%
\setlength\fboxsep{0pt}%
\noindent\protect\colorbox{ocre!40}{\strut\protect\parbox[c][.7cm]{\linewidth}{\Large\sffamily\protect\centering #1\quad\mbox{}}}}{#1}}%
\@endpart}
\def\@endpart{\vfil\newpage
\if@twoside
\if@openright
\null
\thispagestyle{empty}%
\newpage
\fi
\fi
\if@tempswa
\twocolumn
\fi}

%----------------------------------------------------------------------------------------
%	CHAPTER HEADINGS
%----------------------------------------------------------------------------------------

\newcommand{\thechapterimage}{}%
\newcommand{\chapterimage}[1]{\renewcommand{\thechapterimage}{#1}}%
\def\@makechapterhead#1{%
{\parindent \z@ \raggedright \normalfont
\ifnum \c@secnumdepth >\m@ne
\if@mainmatter
\begin{tikzpicture}[remember picture,overlay]
\node at (current page.north west)
{\begin{tikzpicture}[remember picture,overlay]
\node[anchor=north west,inner sep=0pt] at (0,0) {\includegraphics[width=\paperwidth]{\thechapterimage}};
% \draw[anchor=west] (\Gm@lmargin,-9cm) node [line width=2pt,rounded corners=15pt,draw=ocre,fill=white,fill opacity=0.5,inner sep=15pt]{\strut\makebox[22cm]{}};
\draw[anchor=west] (\Gm@lmargin+.3cm,-3cm) node {\huge\sffamily\bfseries\color{black}\thechapter. #1\strut};
% \draw[anchor=west] (\Gm@lmargin+.3cm,-9cm) node {\huge\sffamily\bfseries\color{black}\thechapter. #1\strut};
\end{tikzpicture}};
\end{tikzpicture}
\else
\begin{tikzpicture}[remember picture,overlay]
\node at (current page.north west)
{\begin{tikzpicture}[remember picture,overlay]
\node[anchor=north west,inner sep=0pt] at (0,0) {\includegraphics[width=\paperwidth]{\thechapterimage}};
\draw[anchor=west] (\Gm@lmargin,-9cm) node [line width=2pt,rounded corners=15pt,draw=ocre,fill=white,fill opacity=0.5,inner sep=15pt]{\strut\makebox[22cm]{}};
\draw[anchor=west] (\Gm@lmargin+.3cm,-9cm) node {\huge\sffamily\bfseries\color{black}#1\strut};
\end{tikzpicture}};
\end{tikzpicture}
% \fi\fi\par\vspace*{270\p@}}}
\fi\fi\par\vspace*{70\p@}}}

%-------------------------------------------

\def\@makeschapterhead#1{%
\begin{tikzpicture}[remember picture,overlay]
\node at (current page.north west)
{\begin{tikzpicture}[remember picture,overlay]
\node[anchor=north west,inner sep=0pt] at (0,0) {\includegraphics[width=\paperwidth]{\thechapterimage}};
% \draw[anchor=west] (\Gm@lmargin,-9cm) node [line width=1pt,rounded corners=15pt,draw=ocre,fill=white,fill opacity=0.5,inner sep=15pt]{\strut\makebox[22cm]{}};
\draw[anchor=west] (\Gm@lmargin+.3cm,-3cm) node {\huge\sffamily\bfseries\color{black}#1\strut};
\end{tikzpicture}};
\end{tikzpicture}
% \par\vspace*{270\p@}}
\par\vspace*{70\p@}}
\makeatother

%----------------------------------------------------------------------------------------
%	HYPERLINKS IN THE DOCUMENTS
%----------------------------------------------------------------------------------------

\usepackage{hyperref}
\hypersetup{hidelinks,backref=true,pagebackref=true,hyperindex=true,colorlinks=false,breaklinks=true,urlcolor= ocre,bookmarks=true,bookmarksopen=false,pdftitle={Title},pdfauthor={Author}}
\usepackage{bookmark}
\bookmarksetup{
open,
numbered,
addtohook={%
\ifnum\bookmarkget{level}=0 % chapter
\bookmarksetup{bold}%
\fi
\ifnum\bookmarkget{level}=-1 % part
\bookmarksetup{color=ocre,bold}%
\fi
}
} % Insert the commands.tex file which contains the majority of the structure behind the template


\newcommand{\IGNORE} [1] {#1} %% use to remove comments


\newcommand{\KARTIC} [1] {\IGNORE{\textcolor{red} {KARTIC:#1}}}
\newcommand{\GURPRIT} [1] {\IGNORE{\textcolor{blue} {[GURPRIT:#1]}}}
\newcommand{\WOJCIECH} [1] {\IGNORE{\textcolor{brown} {[WOJCIECH:#1]}}}


% % % % % % % New variables (this course)
\newcommand{\pth} {\ensuremath{\bar{p}}}
\newcommand{\xs} {\ensuremath{j}}
\newcommand{\Ij}[1] {\ensuremath{I_{#1}}}


%%%%%%%%%%%%%%%%%%%%%%%%%%%%%%%%%% OLD STUFF 

% % % % % % %  General
\newcommand{\imag}{\ensuremath{\mathrm{\imath}}}
\newcommand{\MyExp}[1] {\ensuremath{\e^{#1}}}
\newcommand{\cn} [2]{\ensuremath{#1+\imag#2}}
\newcommand{\cp} [2]{\ensuremath{#1 \, \MyExp{-\imag #2}}}
\newcommand{\RV}[1] {\ensuremath{\mathbf{#1}}}
\newcommand{\Real}[1] {\ensuremath{\Re(#1)}}
\newcommand{\Imag}[1] {\ensuremath{\Im (#1)}}
\newcommand{\Exp}[1] {\ensuremath{\langle #1 \rangle}}
\newcommand{\Ampl} [1]{\ensuremath{ |#1 |}}
\newcommand{\Phase} [1]{\ensuremath{ \MyExp{\Phi({#1})}}}
% \newcommand{\Ampl} [1]{\ensuremath{\left| #1 \right|}}
% \newcommand{\Exp}[1] {\ensuremath{\left \langle #1 \right \rangle}}
\newcommand{\Var}[1] {\ensuremath{\mathrm{V} \left(#1\right)}}
\newcommand{\Cov}[2] {\ensuremath{\mathrm{Cov} (#1, #2)}}
\newcommand{\e}{\ensuremath{\mathrm{e}}}
\newcommand{\Z}{\ensuremath{\mathbb{Z}}}
\newcommand{\C}{\ensuremath{\mathbb{C}}}
\newcommand{\R}{\ensuremath{\mathbb{R}}}
\newcommand{\conv}{\ensuremath{\otimes}}
\newcommand{\sinc} [1] {\ensuremath{\mathrm{sinc(#1)}}}
\newcommand{\twopartdef}[4]
{
	\left\{
		\begin{array}{ll}
			#1 & \mbox{if } #2 \\
			#3 & \mbox{} #4
		\end{array}
	\right.
}
\newcommand{\DefInt}[4] {\int\limits_{#1}^{#2} \; #3 \;\; \mathrm{d}#4}
\newcommand{\Int}[2] {\int  #1 \; \mathrm{d}#2}
 
% % % % % % % Variables
\newcommand{\N} {\ensuremath{N}}
\newcommand{\rv} {\ensuremath{\varsigma}}
\newcommand{\Est} {\ensuremath{\mathbb{E}}}
\newcommand{\bias} {\ensuremath{\mathcal{B}}}
\newcommand{\var} {\ensuremath{\mathcal{V}}}
\newcommand{\err} {\ensuremath{\Delta}}
\newcommand{\Estim}[1] {\ensuremath{\RV F_{#1}}}

% % % % % % %  Primal
\newcommand{\sfsym} {\ensuremath{\RV S}}
\newcommand{\sdsym} {\ensuremath{g}}
\newcommand{\ifsym} {\ensuremath{f}}
\newcommand{\x} {\ensuremath{x}}
\newcommand{\w} {\ensuremath{\alpha}}
\newcommand{\xii} {\ensuremath{\RV X _i}}
\newcommand{\wii} {\ensuremath{\w_i}}
\newcommand{\Tx} {\ensuremath{T}}
\newcommand{\sfn} {\ensuremath{\sfsym(\x)}}
\newcommand{\sfnp} {\ensuremath{\sfsym'(\x)}}
\newcommand{\ifn} {\ensuremath{\ifsym(\x)}}
\newcommand{\ifnb} {\ensuremath{\ifsym_{\Pi}(\x)}}
\newcommand{\sdn} {\ensuremath{\sdsym(\x)}}
\newcommand{\boxf} [2]{\ensuremath{\Pi_{#1}^{#2}(\x)}}
\newcommand{\shaf} [1]{\ensuremath{{\bot\hspace{-.7ex}\bot\hspace{-.7ex}\bot}_{#1}(\x)}}
 
% % % % % % %  Fourier
\newcommand{\fv} {\ensuremath{\omega}}
\newcommand{\ifv} {\ensuremath{b}}
\newcommand{\rfv} {\ensuremath{a}}
\newcommand{\amp} {\ensuremath{\rho(\omega)}}
\newcommand{\ph} {\ensuremath{\phi(\omega)}}
\newcommand{\FTsym} [1] {\ensuremath{\hat{#1}}}
\newcommand{\FT} [1] {\ensuremath{\hat{#1}(\fv)}}
\newcommand{\SF} {\ensuremath{\FT{\sfsym}}}
\newcommand{\SFp} {\ensuremath{\FTsym{\sfsym}'(\fv)}}
\newcommand{\IFn} {\ensuremath{\FT{\ifsym}}}
\newcommand{\IFnb} {\ensuremath{\FT{\ifsym_{\Pi}}}}
\newcommand{\SFsym} {\ensuremath{\FTsym{\sfsym}}}
\newcommand{\IFsym} {\ensuremath{\FTsym{\ifsym}}}
\newcommand{\BOXF} [2]{\ensuremath{\FTsym{\Pi}_{#1}^{#2}(\fv)}}
\newcommand{\SHAF} [1]{\ensuremath{\FTsym{{\bot\hspace{-.8ex}\bot\hspace{-.8ex}\bot}}_{#1}(\fv)}}


\begin{document}

%----------------------------------------------------------------------------------------
%	TITLE PAGE
%----------------------------------------------------------------------------------------

\begingroup
\thispagestyle{empty}
\begin{tikzpicture}[remember picture,overlay]
\coordinate [below=12cm] (midpoint) at (current page.north);
\node at (current page.north west)
{\begin{tikzpicture}[remember picture,overlay]
\node[anchor=north west,inner sep=0pt] at (0,0) {\includegraphics[width=\paperwidth]{background}}; % Background image
\draw[anchor=north] (midpoint) node [fill=ocre!30!white,fill opacity=0.6,text opacity=1,inner sep=1cm]
  {\Huge\centering\bfseries\sffamily\parbox[c][][t]{\paperwidth}
   {\centering Fourier Analysis of Sampling Patterns \\ for Rendering: Theory and Practice \\[15pt] % Book title
   {\Large Understanding estimation error in Monte Carlo Image Synthesis}\\[20pt] % Subtitle
   {\huge Kartic Subr, Wojciech Jarosz and Gurprit Singh}}}; % Author name
\end{tikzpicture}};
\end{tikzpicture}
\vfill
\endgroup

%----------------------------------------------------------------------------------------
%	COPYRIGHT PAGE
%----------------------------------------------------------------------------------------

\newpage
~\vfill
\thispagestyle{empty}

\noindent Copyright \copyright\ 2016 XXXX\\ % Copyright notice

\noindent \textsc{Published by Publisher}\\ % Publisher

\noindent \textsc{Best-Course.com}\\ % URL

\noindent Licensed under the Creative Commons Attribution-NonCommercial 3.0 Unported License (the ``License''). You may not use this file except in compliance with the License. You may obtain a copy of the License at \url{http://creativecommons.org/licenses/by-nc/3.0}. Unless required by applicable law or agreed to in writing, software distributed under the License is distributed on an \textsc{``as is'' basis, without warranties or conditions of any kind}, either express or implied. See the License for the specific language governing permissions and limitations under the License.\\ % License information

% \noindent \textit{First printing, March 2013} % Printing/edition date

%----------------------------------------------------------------------------------------
%	ABSTRACT PAGE
%----------------------------------------------------------------------------------------


\pagestyle{plain}
\pagenumbering{roman}

{\baselineskip=18pt

\cleardoublepage
\phantomsection
\addcontentsline{toc}{section}{Abstract}
\section*{Abstract}

Photon density estimation techniques are a popular choice for simulating light transport in scenes with complicated geometry and materials. This class of algorithms can be used to accurately simulate inter-reflections, caustics, color bleeding, scattering in participating media and subsurface scattering. Since its introduction, photon density estimation has been significantly extended in computer graphics with the introduction of: specialized techniques that intelligently modify the positions or bandwidths to reduce visual error using a small number of photons, approaches which eliminate error completely in the limit, and methods that use higher-order samples and queries to reduce error in participating media. This course provides the necessary insight to implement all these latest advances in photon density estimation. The course starts out with a short introduction to photon density estimation using classical photon mapping, but the remainder of the two-part course provides new, hands-on explanations of the latest developments in this area by the experts behind each technique. The course will also cover recent applications of these techniques within industry. The course will give the audience concrete and practical understanding of the latest developments in photon density estimation techniques that have not been presented in prior SIGGRAPH courses.

}

\clearpage


%----------------------------------------------------------------------------------------
%	LECTURERS PAGE
%----------------------------------------------------------------------------------------

\phantomsection
\addcontentsline{toc}{section}{About the Lecturers}
\section*{About the Lecturers}

\subsection*{Prof.\ Kartic Subr}
fixme college\\
\href{mailto:fixme}{\texttt{<fixme.edu>}}\\
\url{http://fixme.com}\\

blah blah blah

\bigskip

\subsection*{Prof.\ Wojciech Jarosz}
Dartmouth College\\
\href{mailto:wojciech.k.jarosz@dartmouth.edu}{\texttt{<wojciech.k.jarosz@dartmouth.edu>}}\\
\url{http://www.cs.dartmouth.edu/~wjarosz}\\

Wojciech Jarosz is an Assistant Professor of Computer Science at Dartmouth College and co-founder of Dartmouth's Visual Computing Lab. Before joining Dartmouth in 2015, Wojciech was a Senior Research Scientist heading the rendering group at Disney Research Zürich, and an adjunct lecturer at ETH Zürich. He obtained his Ph.D. (2008) and M.S. (2005) in computer graphics from UC San Diego, and a B.S. (2003) in computer science from the University of Illinois at Urbana–Champaign.

Prof.\ Jarosz’s research covers many areas of computer graphics and rendering, with a particular emphasis on \emph{light}: how to capture it, how to more efficiently simulate its interaction and propagation within scenes, how to author or edit appearance, and how to create physical objects with prescribed light interactions. Much of his work has involved developing more sophisticated sampling techniques in the context of Monte Carlo integration and the resulting methods have been incorporated into production rendering systems for the making of feature films, including Disney's \textit{Tangled} and \textit{Big Hero 6}. In 2013, Prof.\ Jarosz received the Eurographics Young Researcher Award.

\bigskip

\subsection*{Dr.\ Gurprit Singh}
Dartmouth College\\
\href{mailto:fixme}{\texttt{<fixme.edu>}}\\
\url{http://fixme.com}\\

blah blah blah

\cleardoublepage



%----------------------------------------------------------------------------------------
%	COURSE SYLLABUS
%----------------------------------------------------------------------------------------

\phantomsection
\addcontentsline{toc}{chapter}{Course Syllabus}
\chapterimage{temp.pdf} % Table of contents heading image
\chapter*{Course Syllabus}

Photon mapping techniques have progressed significantly in the last few years. The first part will include a brief introduction to light transport and the core ideas behind photon mapping. The remainder of the time will focus on recent advances in photon density estimation for surfaces.

The second part will start with a brief introduction to participating media, followed by in-depth presentations on recent developments in volumetric photon mapping. The last portion of the course will cover the uses of photon density estimation in industry.

These topics are shown below:

A detailed breakdown in provided in Figure~\ref{fig:half-day}. Each slot is 5~min.

\begin{figure}[b!]
\begin{center}
\end{center}
\caption{%
\label{fig:half-day}%
Session timing breakdown for the course.
}
\end{figure}



%----------------------------------------------------------------------------------------
%	TABLE OF CONTENTS
%----------------------------------------------------------------------------------------

\chapterimage{temp.pdf} % Table of contents heading image

\pagestyle{empty} % No headers

\phantomsection
\addcontentsline{toc}{chapter}{Table of Contents}
\tableofcontents % Print the table of contents itself

\cleardoublepage % Forces the first chapter to start on an odd page so it's on the right

\pagestyle{fancy} % Print headers again

%----------------------------------------------------------------------------------------
%	CHAPTER 1
%----------------------------------------------------------------------------------------

\chapterimage{temp.pdf} % Chapter heading image

\chapter{Introduction}
% Evolution has honed the ability of our visual systems to interpret visual snapshots, or images of our environments, that are sensed by our eyes.
Persuading our sophisticated visual systems of the ``realism'' of computer graphics imagery is a challenging task. 
This problem is exacerbated by the presence of fantastic objects and/or characters embedded in the virtual scene, thereby loosening the notion of realism in the resulting depiction. 
Realistic depiction is typically achieved by an artful combination of design considerations such as the degree of detail in geometrical modelling of objects in the environment, application of plausible models for dynamics and animation of thes objects, the choice of display technology used to present the visuals to the observer, etc. 

While the totality of the above decisions impacts the illusion of reality, the focus of this course is on a vital component in the pipeline --- \textit{the synthesis of photorealistic images} given the various models that describe the environment. The state of the art in the generation of plausible pictures -- which we will henceforth refer to as \textit{rendering} --  involves simulation of the physics of light.
Rendering is a computationally expensive process that needs to be performed repetitively, on every frame, to generate a video sequences.



\section{Rendering}\index{Rendering}
We refer to rendering as the process of producing a picture which represents a particular view of an environment, captured using a virtual camera, at a given time. The input to the rendering system, or renderer, is a detailed representation of all the objects in the scene. This representation is typically provided in a \textit{scene description file} that specifies the objects in the scene, their shapes, their material properties (rough, smooth, shiny, diffuse, translucent, etc.), their motion parameters and whether these objects emit light energy. In addition to the environment, the renderer is also provided with a description of the virtual camera used to capture the rendered view including its position, orientation, optical characteristics (any lenses, aperture, etc.), shutter parameters (exposure time)  and parameters of the virtual sensor (aspect ratio, optical sensitivity, etc.) embedded within it. 

The goal of rendering is to estimate the quantity of radiant light energy impinging on the various pixels (or cells) on the virtual observer's camera (or eye). Current renderers perform this via physically-based simulation of light, following its intricate combinations of reflections within the scene. The virtual measurement \Ij {\xs} at a point \xs\  on the sensor~\cite{VeachChapter8} is
\begin{eqnarray} \label{eq:measeq}
 \Ij {\xs} & = & \DefInt{\Omega}{}{f_\xs(\pth)}{\mu(\pth)}
\end{eqnarray}
where the measurement contribution function, $f_\xs(\pth)$, is the importance-weighted radiance arriving at \xs\ along the path \pth\ and $\mu(\pth)$ is a measure on the space of all paths $\Omega$. The paths \pth\ span space as well as time and may be of arbitrarily high dimensionality. Thus, physically-based rendering requires the esimation of high-dimensional integrals at each point \xs\ on the virtual sensor. 

\TBC

\section{The interplay between reconstruction and integration}\index{Reconstruction and integration}
Images are typically represented and displayed on a regular grid. Rendering algorithms, however, may estimate virtual measurements at arbitrarily points on the sensor. We refer to the process of resampling the measurements \Ij{}\ on a regular grid as \textit{2D reconstruction}. Some rendering algorithms~\cite{MDAS,Egan09,Soleretal09} partition the space of paths ($\Omega$ in equation~\ref{eq:measeq}) explicitly into dimensions that are dependent on the camera, such as exposure time and aperture, and those that depend on the scene. Since the sensor is two-dimensional, exposure time is one-dimensional and the aperture is two-dimensional, they estimate virtual measurements in a five-dimensional, camera-dependent space. This involves a reconstruction in 5D followed by a further projection (integration) down to the two-dimensional plane of the sensor. Such a partitioning is exploited for efficient rendering of camera-dependent blur effects due to fast motion or defocus.

\TBC

\section{The role of sampling}\index{The role of sampling}
There are two distinctly different applications of sampling in the rendering process. One of them is for reconstruction, either in 2D or in 5D (as described above), and the other is for estimation of the integral in equation~\ref{eq:measeq}. Most rendering algorithms focus on the latter and perform reconstruction on a two-dimensional regular grid of pixels. A popular rendering method, called path-tracing, proceeds by estimating the virtual image measurement at each pixel by averaging the values of the measurement contribution function at a chosen set of paths through the pixels. This set of paths may either be sampled stochastically (Monte Carlo sampling) from amongst all paths through each pixel or based on deterministically sampled paths (Quasi-Monte Carlo). The properties of estimates yielded by these two broad classes of algorithms are strikingly different, as is the mathematical machinery used for their analyses.  

\TBC

\section{Sources and manifestation of error}\index{Error in rendering}
The rendered image is represented on a discrete grid, while the distribution of the underlying light energy across the sensor is a continuous function. One potential source of error stems from this discretization. Recall that the true continuous function is a high-dimensional integral at each point on the sensor (equation~\ref{eq:measeq}). Any sampling-based approach used to estimate this function results in an approximation whose measured deviation from the true function is known referred to as \textit{approximation error}. Typically approximation error manifests itself in rendered images in one of two perceivably objectionable forms: structured artifacts or as noise. Various metrics, such as L1 norms, L2 norms, PSNR, etc., are used to quantify approximation error.

\TBC 
%----------------------------------------------------------------------------------------
%	CHAPTER 2
%----------------------------------------------------------------------------------------
\chapter{Mathematical preliminaries}
\section{Random variables, expectation and variance}
\section{Fourier series}
\section{The continuous Fourier transform}
\section{The Discrete Fourier transform}
\section{The Fast Fourier transform}

%----------------------------------------------------------------------------------------
%	CHAPTER 3
%----------------------------------------------------------------------------------------
\chapter{Image reconstruction}

\section{The visual system and perceptual reconstruction}
\section{Aliasing and error}
\section{Quantitative assessment}
\section{Antialiasing}

%----------------------------------------------------------------------------------------
%	CHAPTER 4
%----------------------------------------------------------------------------------------
\chapter{Estimating integrals}

\section{Numerical integration via sampling}
\section{Error and convergence}
\section{Error in the Fourier domain due to sampling}
\section{Uncertainty of estimates}
\section{Assessing sampling patterns based on their spectra}
\section{Analysis beyond the canonical domain}
\subsection{Spherical domain}
\subsection{General domain}
\subsection{Gradient domain}

%----------------------------------------------------------------------------------------
%	CHAPTER 5
%----------------------------------------------------------------------------------------
\chapter{Popular sampling strategies}

\section{Classical}
\subsection{Random sampling}
\subsection{Jittered sampling}
\subsection{Uniform jittered sampling}

\section{Quasi-Monte Carlo}
\subsection{Sobol sequence}
\subsection{Hammersley sequence}
\subsection{Latin hypercube sampling}

\section{Blue noise}
\subsection{Poisson-disk sampling}
\subsection{CCVT sampling}
\subsection{Tiling-based methods}

\section{Synthesis of sampling patterns with targeted spectral profiles}

%----------------------------------------------------------------------------------------
%	CHAPTER 6
%----------------------------------------------------------------------------------------
\chapter{Case studies}
\section{1D integrands}
\subsection{Smooth functions}
\subsection{Functions with discontinuities}
\section{2D integrands}
\subsection{Smooth functions}
\subsection{Sub-pixel integration}
\subsection{Functions with discontinuities}
\section{Rendering test suite}

% 
% \section{Figure}\index{Figure}
% 
% \begin{figure}[h]
% \centering\includegraphics[scale=0.5]{placeholder}
% \caption{Figure caption}
% \end{figure}

%----------------------------------------------------------------------------------------
%	BIBLIOGRAPHY
%----------------------------------------------------------------------------------------

\chapter*{Bibliography}
\addcontentsline{toc}{chapter}{\textcolor{ocre}{Bibliography}}
\section*{Books}
\addcontentsline{toc}{section}{Books}
\printbibliography[heading=bibempty,type=book]
\section*{Articles}
\addcontentsline{toc}{section}{Articles}
\printbibliography[heading=bibempty,type=article]

%----------------------------------------------------------------------------------------
%	INDEX
%----------------------------------------------------------------------------------------

\cleardoublepage
\phantomsection
\setlength{\columnsep}{0.75cm}
\addcontentsline{toc}{chapter}{\textcolor{ocre}{Index}}
\printindex

%----------------------------------------------------------------------------------------

\end{document}